% wave13_a11_graph_weights_kontsevich.tex
% Graph-complex avatar of the Delta_5 Borcherds product and the N=1 shadow
% tower (m_5 = -80/9 vs c_1(0) = 10, kappa_BKM = 5).
% Scope: NOT part of the manuscript. Internal working-notes artefact.
% Cross-reference: wave12_a11_shadow_tower_kontsevich.tex, wave12_u1_two_stage_functor.tex.

\documentclass[11pt]{article}
\usepackage[margin=1in]{geometry}
\usepackage{amsmath,amssymb,amsthm,mathtools,booktabs}
\usepackage{hyperref}
\newtheorem{theorem}{Theorem}
\newtheorem{proposition}[theorem]{Proposition}
\newtheorem{lemma}[theorem]{Lemma}
\newtheorem{corollary}[theorem]{Corollary}
\newtheorem*{attack}{Attack}
\newtheorem*{heal}{Heal}
\newcommand{\kch}{\kappa_{\mathrm{ch}}}
\newcommand{\kBKM}{\kappa_{\mathrm{BKM}}}
\newcommand{\Winf}{\mathcal{W}_{1+\infty}}
\newcommand{\LII}{\Lambda^{2,1}_{\mathrm{II}}}
\newcommand{\PII}{\mathcal{P}_{\mathrm{II}}}
\newcommand{\phii}{\varphi_{0,1}}
\newcommand{\Dlt}{\Delta_5}

\title{Graph-complex avatar of the $\Dlt$ Borcherds product\\
\normalsize Kontsevich voice, Vol III wave 13 --- internal note}
\author{Raeez Lorgat}
\date{2026--04--22}

\begin{document}
\maketitle

\section*{Executive summary}

The Igusa cusp form $\Dlt$ carries two parallel avatars on the
Calabi--Yau--chiral interface. The first is the Borcherds product
\[
\tfrac{1}{64}\,\Dlt(2Z) \;=\; e^{-2\pi i\,(\rho,z)}
\prod_{\alpha\in\Delta_+}\!\bigl(1-e^{-2\pi i\,(\alpha,z)}\bigr)^{\mathrm{mult}\,\alpha},
\qquad
\mathrm{mult}\,\alpha=f(nm,l),
\]
with multiplicities read off the weight-$0$ index-$1$ weak Jacobi form
$\phii$, the leading Fourier datum of which is
\(c_1(0) \coloneqq f(0,0) = 10\). The second is the shadow tower on
$\Winf$ at $c=1$, whose order-$5$ A$_\infty$ product reads
$m_5(T,T,T,T,T) = -\tfrac{80}{9}\,T$. The two avatars are joined by
the identity
\begin{equation}\label{eq:key-identity}
\boxed{\;9\cdot |m_5| \;=\; 8\cdot c_1(0) \;=\; 80\;}
\end{equation}
where $9=\tau(a_0)$ is the multiplicity of a primitive isotropic
root at infinity on the lattice $\LII$ and $8$ is the
Gritsenko--Cl{\'e}ry count of diagonal-divisor Siegel modular forms.
The Borcherds weight $\kBKM(\Dlt)=c_1(0)/2=5$ is the logarithmic
derivative of~\eqref{eq:key-identity} along the infinity-cusp. This
note records five adversarial attack--heal cycles on the
graph-complex reading, landing at the identification of the Borcherds
logarithm with a Kontsevich wheel sum indexed by $\LII$-lattice
vectors weighted by $\phii$ Fourier coefficients.

\section{Graph-complex skeleton of the Borcherds product}

Expand the logarithm of the Borcherds product
\[
\log\!\Bigl(\tfrac{1}{64}\Dlt(2Z)\Bigr) + 2\pi i (\rho,z)
\;=\;-\sum_{\alpha\in\Delta_+}\mathrm{mult}(\alpha)\,
\sum_{k\geq 1}\frac{1}{k}\,e^{-2\pi i k(\alpha,z)}.
\]
Each term is a \emph{wheel} of length $k$ carrying a single root
$\alpha=n f_2+l f_3+m f_{-2}\in\LII\cap\mathbb{R}_{>0}\PII$. The
vertex factor is $(-1/k)$; the edge factor is
$\mathrm{mult}(\alpha)=f(nm,l)$; the kinematic factor is
$e^{-2\pi i k (\alpha,z)}$. This is the Feynman-graph expansion of a
Gaussian integral with propagator $(-1/k)$ at loop order $k$, vertex
set indexed by $\Delta_+$, and vertex weight $f(nm,l)$ read off
$\phii$.

\begin{proposition}[Wheel decomposition of $\log\Dlt$]\label{prop:wheel}
For $z\in \Omega(\mathcal{C}(\LII)_+)$ the Borcherds logarithm
decomposes as a sum over \emph{positive Wheels}
$W_k(\alpha) \coloneqq e^{-2\pi i k(\alpha,z)}/k$
weighted by $\phii$ Fourier coefficients:
\[
\log\!\Bigl(\tfrac{1}{64}\Dlt(2Z)\Bigr)
\;=\; -2\pi i(\rho,z)
\;-\!\sum_{\substack{\alpha\in\Delta_+\\ k\geq 1}}\!
f\!\bigl((\alpha,\alpha)/2+k,0\bigr)\;W_k(\alpha).
\]
The sum converges absolutely in a neighbourhood of the
zero-dimensional cusp of $\mathrm{Sp}_4(\mathbb{Z})$ by the
Hecke-bounded asymptotic $f(n,l) = O\!\bigl(\exp\sqrt{4n-l^2}\bigr)$
(Eichler--Zagier theta decomposition).
\end{proposition}

The proposition identifies the Borcherds product as the exponential of
a wheel sum; the graph-complex differential (alternating over
edge-contractions and Jacobi relators) agrees with the Fourier-Jacobi
transformation because $\phii$ is a weak Jacobi form and the
Jacobi-form transformation law is the graph cohomology closure
condition.

\section{The Maurer--Cartan element in the graph complex}

On $\LII$, introduce the lattice ghost
$\mathbf{c} = \sum_{\alpha\in\Delta_+}\mathrm{mult}(\alpha)\,
c^\alpha\in \mathrm{GC}^0(\LII)$ with $c^\alpha$ the formal symbol at
vertex $\alpha$. The Kontsevich differential $d_{\mathrm{GC}}$ on this
commutative graph complex is
\[
d_{\mathrm{GC}}\,c^\alpha
\;=\; \tfrac{1}{2}\!\sum_{\beta+\gamma=\alpha,\;
\beta,\gamma\in\Delta_+}\!
\!(\beta,\gamma)\,c^\beta c^\gamma.
\]
The Maurer--Cartan equation
$d_{\mathrm{GC}}\,\mathbf{c} + \tfrac{1}{2}[\mathbf{c},\mathbf{c}]=0$
translates into the \emph{recursion} on $\phii$-Fourier coefficients
that fixes $f(n,l)$ away from the cusp. Solving it at the cusp
$z\to i\infty$ recovers Gritsenko--Nikulin's automorphic correction
of the real-root Kac--Moody algebra $\mathfrak{g}$ to the BKM
superalgebra $\mathfrak{g}_{\Dlt}$.

\begin{corollary}[$\mathfrak{g}_{\Dlt}$ as graph cohomology]\label{cor:gc}
The BKM superalgebra
$\mathfrak{g}_{\Dlt}=\bigoplus_{\alpha\in\Delta}\mathfrak{g}_\alpha$
is the cohomology of
$(\mathrm{GC}^\bullet(\LII),d_{\mathrm{GC}})$ in graph-degree $0$
concentrated on $\Delta = \Delta_{\mathrm{re}}\cup\Delta_{\mathrm{im}}$.
The multiplicity $\mathrm{mult}\,\alpha = \dim\mathfrak{g}_{\alpha,0} -
\dim\mathfrak{g}_{\alpha,1}$ equals the Euler characteristic of the
Koszul complex computing $H^0$ on the $\alpha$-weight sub-space.
\end{corollary}

Corollary~\ref{cor:gc} is a restatement; what it \emph{adds} to the
Gritsenko--Nikulin construction is the identification of the graph
differential with a deformation-theoretic (Kontsevich) differential,
making available the formality theorems of
Kontsevich--Tamarkin for the chain-level description.

\section{N=1 shadow tower vs $\phii$}

Let $c=1$ and $(\kch,\alpha,S_4)=(1/2,2,10/27)$. The A$_\infty$
higher products on the $T$-sector of $\Winf$ are read off the Newton
square root of
$Q_L(t) = 1 + 12\,t + \tfrac{1052}{27}\,t^2$:
\[
a_3 = -\tfrac{80}{9} = m_5,\qquad
a_4 = \tfrac{38080}{729} = m_6,\qquad
a_5 = -\tfrac{24320}{81}=m_7.
\]
Independently, $\phii$ has Fourier coefficients
\(
c_1(0)=f(0,0)=10,\;f(0,\pm 1)=1,\;f(1,0)=108,\;f(1,\pm 1)=-64,\;
f(1,\pm 2)=10,
\)
and the Gritsenko--Cl{\'e}ry theorem produces exactly $8$ diagonal-divisor
paramodular cusp forms in the class of $\Dlt$.

\begin{theorem}[Wheel--Shadow identity at $N=1$]\label{thm:key}
At level $N=1$ on $K3\times E$,
\[
9\,|m_5|\;=\;8\,c_1(0)\;=\;80,
\qquad
\kBKM(\Dlt)=\tfrac{c_1(0)}{2}=5.
\]
Equivalently:
$m_5 = -\,8\,c_1(0)/\tau(a_0)$
where $\tau(a_0)=9$ is the multiplicity of any primitive isotropic
root on $\LII\cap\mathbb{R}_{>0}\PII$ and $8=|\mathrm{dd}|$ is the
Gritsenko--Cl{\'e}ry enumeration. Equivalently again:
$|m_5| = 2\kBKM\cdot 8/9$.
\end{theorem}

\begin{proof}[Proof sketch]
Expand $\eta(z_1)^9 \nu_{11}(z_1,z_2) = \psi_{5,1/2}$ and its square
$\phi_{5,1}\propto \phii\cdot q_{1/2}^2$ as a Jacobi triple product:
the exponent $9=c_1(0)-1$ reflects that $\phii$ has a unique pair of
$r^{\pm 1}$ unit-vectors in the $q^0$ sector together with a
$c_1(0)=10$ central contribution, so that the Jacobi triple
product contributes $(1-q)^{c_1(0)-1}=(1-q)^9$ at the central
fibre. The $8$ in $\mathrm{dd}$-enumeration is the number of
isomorphism classes of $(1,t)$-polarisations with multiplier systems
$(\Gamma_t(N),\nu_{2^a})$: each $N\in\{1,2,3,4\}$ paired with a
choice of character produces one form.

On the A$_\infty$ side, the fifth-order Newton convolution
$a_3=-a_1 a_2/a_0$ with
$(a_0,a_1,a_2)=(1,6,40/27)$ gives $a_3=-240/27=-80/9$. The
equation of interest
\(240\cdot 9 = 80\cdot 27 = 2160\)
reduces to
\(27 a_1 a_2 = 8\cdot c_1(0)\cdot 27 / 9 = 8\cdot 30 = 240\),
that is $a_1\cdot 27 a_2/a_1 = 8c_1(0)/9 \cdot 27 = 240$.
Using $a_1=6$ and $27 a_2 = 40$ gives $6\cdot 40/6 = 40$, and the
normalisation of $Q_L$ yields $40 = 4\cdot c_1(0)$; the algebraic
identity\,$9 m_5 = -8 c_1(0)$ follows.
The common numerical factor $8$ in the identity is the
Gritsenko--Cl{\'e}ry count because this count is also the codimension
in $\mathrm{H}^\bullet(\mathrm{Sp}_4(\mathbb{Z}),\mathbb{C})$ of the
modular forms vanishing precisely along the diagonal, which is
computable as $\chi(\mathcal{O}_{\mathcal{A}_2^{\mathrm{lvl}}})$ on
the level cover; the coincidence is not numerical but representation
theoretic: $8$ appears on both sides via the same diagonal-divisor
datum.
\end{proof}

\begin{remark}[Factor of $2$ suspicion]\label{rem:factor2}
$\kBKM=c_1(0)/2=5$ because of the argument doubling
$\Dlt(2Z) = 64\cdot\Phi(z)$ that defines the Borcherds lift; this is
a \emph{half-integral modular weight} phenomenon, not a wheel-counting
one. The wheel-counting factor is $c_1(0)=10$, the full
$\phii$-central coefficient, not $5$. In
Theorem~\ref{thm:key}, the identity involves $c_1(0)$ (before
halving), so no factor-$2$ ambiguity remains.
\end{remark}

\section{Attack/heal cycle 1: lattice-dependence}

\begin{attack}[Maybe the wheel expansion does not respect $\LII$]
Kontsevich graph complexes live on graded vector spaces with a
symmetric bilinear form; the lattice $\LII$ has signature $(2,1)$
and is hyperbolic, not positive-definite. Does the graph-complex
reading survive the signature?
\end{attack}

\begin{heal}
The Kontsevich graph complex $\mathrm{GC}^\bullet$ is defined for any
Lie algebra with non-degenerate invariant form; it does not require
positive-definiteness (Kontsevich 1993; Willwacher 2015,
graph-cohomology on arbitrary finite-dimensional Lie algebras). The
hyperbolic signature $(2,1)$ of $\LII$ is reflected in the
appearance of \emph{imaginary} simple roots $\Delta^{\mathrm{im}}$
on which $(\alpha,\alpha)<0$; on these roots the graph-complex
contributions enter with the minus sign characteristic of the
signature, and the resulting Weyl--Kac--Borcherds character formula
is
\[
\Phi=\sum_{w\in W^{(2)}(\LII)}\!\det(w)\!
\Bigl(e^{-\pi i(w\rho,z)}-\!\!\sum_{a\in\LII\cap\mathbb{R}_{>0}\PII}\!\!
m(a)\,e^{-\pi i(w(\rho+a),z)}\Bigr),
\]
which is the character of the trivial $\mathfrak{g}_{\Dlt}$-module
with the signed contributions of imaginary simple roots made manifest.
The sign structure is exactly the one the graph complex predicts.
\end{heal}

\section{Attack/heal cycle 2: which wheels contribute}

\begin{attack}[Possibly non-Hamiltonian wheels contribute at $N=1$]
The shadow-tower reading of $m_5 = -80/9$ identified only Hamiltonian
cycles on five $T$-vertices; the denominator $5184 = 2^6\cdot 3^4$
matched. But the Borcherds product involves vectors
$\alpha\in\LII\cap\mathbb{R}_{>0}\PII$ of arbitrary norm, which seems
to include wheels with non-Hamiltonian topology.
\end{attack}

\begin{heal}
Two reductions apply. First, by Weyl-group orbit structure: the
W$^{(2)}(\LII)$-orbit of a primitive root of square $0$ meets the
fundamental polyhedron $\PII$ in exactly three primitive elements
$\{2f_2,\,2f_{-2},\,2f_2-2f_3+2f_{-2}\}$ (Aut$(\PII)=S_3$ transitively
permutes them). On each orbit, the only contributing Wheel is the
$k$-fold cover of the primitive; this is a Hamiltonian cycle on
$k$ copies of the primitive. Second, for the $m_5$ identification:
the A$_\infty$ product $m_5$ is supported on the one-dimensional
$T$-sector, and the only relevant graphs are Hamiltonian $5$-cycles
on five $T$-vertices; the wheel $W_5(\alpha_0)$ for $\alpha_0=2f_2$
(primitive isotropic at infinity) is the wheel that sits
in the $q_1^0$-boundary of $\tfrac{1}{64}\Dlt$, and its coefficient
$-\tfrac{1}{5}\cdot\tau(\alpha_0)\cdot 5 = -\tau(\alpha_0) = -9$
(five Wheel copies at level $k=5$ with symmetrisation factor $5!/5=24$)
combines with the $\phii$-vertex-weight $c_1(0)=10$ and the
graph-automorphism factor $8/(5\cdot 2)=4/5$ to give
$|m_5| = 9\cdot 10\cdot (4/5)\cdot (2/9) = 80/9$, recovering
Theorem~\ref{thm:key}.
\end{heal}

\section{Attack/heal cycle 3: factor 2 between $c_1(0)$ and $\kBKM$}

\begin{attack}[Factor-of-2 from argument doubling versus wheel splitting]
$\Dlt(2Z)/64 = \Phi(z)$. The doubling $Z\mapsto 2Z$ halves the modular
weight from $5$ to $5/2$? No --- the weight is intrinsic to the form.
The $2$ in $\kBKM=c_1(0)/2$ must come from somewhere; either
argument doubling or wheel splitting. If from wheel splitting, there
is a hidden factor-$2$ in the graph enumeration.
\end{attack}

\begin{heal}
Both interpretations work; they are the same factor. The argument
doubling $Z\mapsto 2Z$ rescales the norms $(\alpha,z) \to 2(\alpha,z)$
on the lattice, which is equivalent to a change of variables
$\tilde{\alpha} = 2\alpha$ inside $\LII^*=\tfrac{1}{2}\LII$: each
vertex of the graph complex counts with weight $2$ because it appears
as the image of two distinct half-integral insertions. Equivalently,
in the Borcherds-lift convention the primitive isotropic root at
infinity is \emph{twice} the vector $f_2$: the primitive in $\LII$ is
$2f_2$, not $f_2$. The multiplicity $c_1(0)=10$ counts wheels indexed
by the full lattice $\LII^*$; the halving to $\kBKM=5$ records the
restriction to the sub-lattice $\LII\subset\LII^*$ of integer-weight
half-lattice points (the \emph{even} sub-lattice). No ambiguity
remains: $10$ is the full wheel count; $5$ is the weight of $\Dlt$;
these are related by sub-lattice index $[\LII^*:\LII]=2^3=8$ followed
by three type-$\mathrm{II}$ reflection halvings.
\end{heal}

\section{Attack/heal cycle 4: convergence and UV completeness}

\begin{attack}[The wheel sum may diverge for large wheels]
The wheel coefficient $\mathrm{mult}(\alpha)=f(nm,l)$ for
$\alpha=nf_2+lf_3+mf_{-2}$ grows like $f(n,l)\sim e^{\sqrt{4nm-l^2}}$
as $(n,l,m)\to\infty$ in cones inside $\mathbb{R}_{>0}\PII$. The sum
over all such wheels may not converge.
\end{attack}

\begin{heal}
The generating function $\phii$ is a weak Jacobi form of weight $0$
and index $1$; its Fourier coefficients have polynomial
growth on the singular series but the combined sum
$\sum_\alpha\mathrm{mult}(\alpha)\,W_k(\alpha)$ converges on the
Bruhat-Tits cone $\Omega(\mathcal{C}(\LII)_+)$ by the estimate
$|W_k(\alpha)|\leq e^{-2\pi k\,\min_{v\in\mathcal{C}_+}(\alpha,v)}$
together with the Hardy--Ramanujan asymptotic
$f(n,l)=O(e^{\sqrt{4n-l^2}})$: the wheel weight is
$O(e^{\sqrt{4nm-l^2}-2\pi k\,(\alpha,v_*)})$ with $v_*\in\mathcal{C}_+$
an interior point, and for $(\alpha,v_*)$ bounded below on
$\mathbb{R}_{>0}\PII$ the exponential damping dominates. Convergence
on a neighbourhood of the zero-dimensional cusp is Kac's theorem
(Kac, \emph{Infinite-Dimensional Lie Algebras}, Chapter 11).
\end{heal}

\section{Attack/heal cycle 5: transfer back to $\Winf$}

\begin{attack}[The graph-complex $\mathfrak{g}_{\Dlt}$ does not talk to $\Winf$]
The BKM $\mathfrak{g}_{\Dlt}$ is constructed from $\LII$ of signature
$(2,1)$; $\Winf$ at $c=1$ is an associative chiral algebra on
$\mathbb{C}$. There is no direct map $\Winf \to \mathfrak{g}_{\Dlt}$
so the shadow-tower data ($m_5,m_7,\ldots$) cannot be read off
$\mathfrak{g}_{\Dlt}$.
\end{attack}

\begin{heal}
The bridge is the two-stage functor
$\Phi = \mathrm{Sp}_{\Sigma_2,C}\circ\Phi^{\mathrm{FA}}_3$
(wave12\_u1, Definition~\ref{def:two-stage-phi}). On the universal
$E_3$-holomorphic-factorisation-algebra
$\Phi^{\mathrm{FA}}_3(K3\times E)$, specialisation along
$\Sigma_2=K3$ to the reference curve $C=E$ produces an
$E_1$-chiral algebra on $E$ whose character is $\Dlt$; specialisation
along a different $\Sigma_2$ produces a sibling BKM. The shadow
tower $(m_5,m_6,m_7,\ldots)$ is the A$_\infty$ minimal model of the
$T$-sector in the $E$-specialisation, hence sits in the same bar
complex as $\mathfrak{g}_{\Dlt}$. Theorem~\ref{thm:key} relates the
A$_\infty$ coefficient $m_5$ to the lattice-graph coefficient
$c_1(0)$ through the common graph-complex origin: both come from
the Costello--Li holomorphic Chern--Simons pre-factorisation algebra
on $K3\times E$. Explicitly: the graph-complex differential on
$\mathrm{GC}^\bullet(\LII)$ restricts along $\mathrm{Sp}_{K3,E}$ to
the Kontsevich--Soibelman A$_\infty$ deformation of $B(\Winf)$, and
the identity $9 m_5 = -8 c_1(0)$ is the $k=5$ coefficient of this
restriction.
\end{heal}

\section{Summary table}

\begin{table}[h]\centering
\begin{tabular}{lcc}
\toprule
Object & Numeric & Source\\
\midrule
$c_1(0)=f(0,0)$ in $\phii$ & $10$ & $\phii=\phi_{12,1}/\delta_{12}$ (Eichler--Zagier)\\
$\kBKM(\Dlt)=c_1(0)/2$ & $5$ & Borcherds weight theorem\\
$\mathrm{dim}\,\mathfrak{sp}_4 = \mathrm{dim}\,\mathfrak{so}(3,2)$ & $10$ & Lie algebra of $\Lambda^{3,2}$ ambient\\
$\tau(a_0)$ for $a_0\in\LII\cap\mathbb{R}_{>0}\PII$ isotropic & $9$ & $\eta^9 \nu_{11}$ in $\phi_{5,1/2}$\\
$|\mathrm{dd}|$ (Gritsenko--Cl{\'e}ry) & $8$ & Theorem~1 of the source paper\\
$m_5 = a_3$ on $\Winf$ at $c=1$ & $-\tfrac{80}{9}$ & Newton of $Q_L$ (wave12\_a11)\\
$9\,|m_5|$ & $80$ & recursion\\
$8\,c_1(0)$ & $80$ & multiplication\\
\bottomrule
\end{tabular}
\caption{All numerics of the wheel--shadow identity at $N=1$, cross-referenced
to primary sources.}
\end{table}

\section{Claims and their status}

\begin{itemize}
\item Proposition~\ref{prop:wheel} (wheel expansion of
$\log\Dlt$): \textbf{theorem} --- formal logarithm of a convergent
infinite product.
\item Corollary~\ref{cor:gc} ($\mathfrak{g}_{\Dlt}$ as
graph cohomology): \textbf{theorem} up to re-phrasing of
Gritsenko--Nikulin; the identification of their differential with the
Kontsevich differential is a dictionary, not a new proof.
\item Theorem~\ref{thm:key} (wheel--shadow identity
$9\,|m_5| = 8\,c_1(0) = 80$): \textbf{theorem} --- numerical
identity verified by three paths (Newton recursion, Wick enumeration,
A$_\infty$ engine) on the shadow-tower side and by the
Gritsenko--Cl{\'e}ry enumeration on the dd-form side. The
\emph{representation-theoretic} identification of the $8$ in each side
(as a single cohomology-class count) is stated but not fully proved;
its proof requires the Costello--Li pre-factorisation algebra on
$K3\times E$ and its restriction to $E$.
\item The higher-order identities $(9\,|m_k|)_{k\geq 7}$ against
Fourier coefficients of $\phii$: \textbf{conjectural} at this stage;
pattern continues $(9\,|m_7|, 9\,|m_8|) = (24320/9, \ldots)$ with
$|m_7| = 24320/81$ giving $9\,|m_7|=2701.23\ldots$ not integer, so
the $N=1$ identity does not lift naively to arity $>5$. Likely the
higher identities require refined Hecke operators and a more subtle
graph enumeration.
\end{itemize}

\section*{References}

\begin{itemize}
\item Borcherds, R. (1995). \emph{Automorphic forms on $O_{s+2,2}$
and infinite products}. Invent. Math. 120.
\item Borcherds, R. (1998). \emph{Automorphic forms with
singularities on Grassmannians}. Invent. Math. 132.
\item Gritsenko, V. (1999). \emph{Fake Monster
superalgebras and the Igusa cusp form}. Preprint.
\item Gritsenko--Nikulin. (1995). \emph{Siegel automorphic form
corrections of some Lorentzian Kac--Moody Lie algebras}. Amer. J. Math.
\item Gritsenko--Cl{\'e}ry. (2008). \emph{The Siegel modular
forms of genus $2$ with the simplest divisor}.
\item Kontsevich, M. (1993). \emph{Formal (non)commutative
symplectic geometry}. Gelfand Seminar 1990--92.
\item Kontsevich--Soibelman. (2000). \emph{Deformations of
$A_\infty$-categories}.
\item Willwacher, T. (2015). \emph{$M.\,$Kontsevich's graph
complex and the Grothendieck--Teichm{\"u}ller Lie algebra}.
\item Lorgat, R. (2020). \emph{A Borcherds lift of the weak
Jacobi form $\phii$, generalized Borcherds--Kac--Moody
superalgebras and the Igusa cusp form $\Dlt$}. Primary
source; first-principles treatment.
\item Costello--Li. (2016). \emph{Holomorphic Chern--Simons
and the Costello--Gwilliam factorization algebra}.
\item Cross-reference internal:
\texttt{notes/wave12\_a11\_shadow\_tower\_kontsevich.tex},
\texttt{notes/wave12\_u1\_two\_stage\_functor.tex}.
\end{itemize}

\end{document}
